\chapter{TPI: Introduction}

	\noindent
	Here begins the second mathematically impossible point. It is a large travel, but we did the half of it in inverse way. I have showed you how to solve it, or how to build it for $\mathcal{P(\mathbb{N})}$ vs $\mathbb{N}$. AND AFTER, I am going to define what we were doing. Normally I do it in an inverse way: first I define TPI, and AFTER I apply it, once people have said me "I dare you to apply it for $\mathcal{P(\mathbb{N})}$ vs $\mathbb{N}$ ".
	\\\\
	
	\noindent
	TPI (Unlimited Transference of Pairs), when it is explained for finite spectrum, or known cases... people accept it, more quickly or more slowly. But they often accept it. After seeing it can be build for cases they never imagined before, they regret their words, very dishonestly...
	\\\\
	
	\noindent
	We are trying to build a proof with two paths. No matter if you try to deny some point of TPI: it has CLEAR consequences doing team with DI. Denying it as a correct tool to compare infinite cardinals is... embarassing. When you saw it in more simple cases and you had judged it as "valid"(I know... I am talking about your recent future :D)
	\\\\
	
	\noindent
	In DI, we use D-pairs, to try to left empty (or NONE empty) a subset of $\mathbb{N}$: an arbitrary MEGAPACK. In TPI, we will use isolated universes, to try to left empty (or not empty) a set of D-pairs: SNEIs X SNEIs. The totallity of it.
	\\\\
	
	\noindent
	We will invert the roles: now universes will help us to discard Families of D-pairs. And we will find another point of view for the "shit-tuation": now we will need S totally empty, but S is going to be (SNEIs X SNEIs), instead the MEGAPACK.
	\\\\
	
	\noindent
	I have a document of 60 pages defining the TPI. JUST defining it. And trying to defend WHY it is a valid tool to compare sets with infinite cardinals:\\\\
	https://vixra.org/abs/2209.0120 \\\\
	It is in spanish. I won't put here all the arguments, and NOW I have the system of relations builded, I don't need to do so much vague explanations. We can do this, "more quick", if you just admit three simple concepts:\\\\
	1.- To solve a certain quantity of D-pairs, we need to admit the Image set MUST have "enough variety" (quantity) of elements to being able to do that.\\\\
	2.- If solving all D-pairs means having, at least, the same "cardinal" as the Domain, we can measure indirectly, improvements, in proven "quantities", after splitted ALL (SNEIs X SNEIs) in $\aleph_{0}$ steps/pieces (Families), and seeing how each universe is able to solve more and more Families. It is an animation in $\aleph_{0}$ steps, from an undeterminated "quantity" of elements, to the quantity of all elements inside SNEIs, but measured indirectly by the quantity of D-pairs we can solve. With each $\theta_{k}$ universe, in each $r_{\theta_{k}}$ relation.\\\\
	3.- The semi-xmas-tree pattern, that in finite spectrum means we have much more elements inside the ORIGIN set, than inside the Domain set, CAN NOT CHANGE its interpretation in infinite spectrum, to the total opposite. We can change it to "equal quantities", but not to "unimaginably smaller".  It is worst than saying doubling the elements inside a set with infinite cardinal, can create a set with a smaller aleph as cardinal. THE MAGNITUDE OF THAT ABSURD IS BEYOND MY HABILITIES AS POET.
	\\\\
	
	\noindent
	We will loose the GREAT MOMENT when someone says: "It is mathematically impossible", to build a TPI for $\mathcal{P(\mathbb{N})}$ (Domain) vs $\mathbb{N}$ (Origin), because I have showed you the system of relations, but very probably we won't loose the moment you will see yourself trying to defend an option for the "Shit-tuation", you considered before as a typical "crankery" answer :D.
	\\\\
	
	\noindent
	Instead of using FLJA\_abstract relation, we will see each $r_{\theta_{k}}$ relation as individuals, talking about simple infinite PACKS.
	\\\\
	
	\noindent
	I have mentioned many concepts... and we need more... Don't worry if you consider the future examples as "insultingly simple". THAT IS THE FRUSTATING POINT: the observation of the properties is very simple. Until the point you will probably shortcircuit trying to deny them, because they are so simple, that are very hard to deny. This is not going to be short. This is going to be large, but not complicated.
	\\\\
	
	\noindent
	Using simple tools to create solutions with huge implications, normally is called "The beauty in mathematics"... but only if you are famous. I am happy if you just consider them as "correct", or "worthy of a deeper study". First angry reaction is assumed, but I rather honesty. So take your time.
	\\\\
	
	\noindent
	1.- We will explain the concept of TPI (solving D-pairs as indirect measure, NWSP, WSP, QRE and semi-xmas-tree pattern) for finite quantities of elements.\\\\
	2.- We will translate it to infinite quantities, because between two consecutive alephs, we don't have different cardinals.\\\\
	3.- WE will see what we have in common with the Shit-tuation.\\\\
	4.- We will discuss the two possible options.\\\\
	5.- Final conclussion merging DI and TPI: a proof with two paths driving to the same conclussion. Third mathematically impossible point: we don't need a proper injectivity.\\\\
	
	
	
		
	
	
	
